% Optional checked bibliography blocks. Merge into the existing bibliography;
% do not include this file wholesale if any key is already present.
% Statement extraction status and scope are recorded in CHECKED_SOURCES.md.

\bibitem{DeligneHodgeStatement}
P. Deligne.
\newblock The Hodge Conjecture.
\newblock Clay Mathematics Institute, official problem description, \S1--2.
\newblock \url{https://www.claymath.org/wp-content/uploads/2022/06/hodge.pdf}.

\bibitem{DeligneHodgeII}
P. Deligne.
\newblock Th\'eorie de Hodge II.
\newblock \emph{Publications Math\'ematiques de l'IH\'ES}, 40:5--57, 1971.
\newblock doi:10.1007/BF02684692.
\newblock Corollary 3.2.17, p.~39; Lemma 4.2.3, p.~44.

\bibitem{ArapuraUniruled}
D. Arapura.
\newblock Varieties with very little transcendental cohomology.
\newblock arXiv:0706.2506v1, 17 June 2007, Corollary 4.2.

\bibitem{ConteMurre}
A. Conte and J. P. Murre.
\newblock The Hodge conjecture for fourfolds admitting a covering by rational curves.
\newblock \emph{Mathematische Annalen}, 238:79--88, 1978.
% Original not freshly extracted; the checked primary reproof is ArapuraUniruled.

\bibitem{BDPP}
S. Boucksom, J.-P. Demailly, M. P\u{a}un, and T. Peternell.
\newblock The pseudo-effective cone of a compact K\"ahler manifold and varieties of negative Kodaira dimension.
\newblock arXiv:math/0405285v1, 14 May 2004, Theorem 2.6.

\bibitem{Aoki}
N. Aoki.
\newblock Some new algebraic cycles on Fermat varieties.
\newblock \emph{Journal of the Mathematical Society of Japan}, 39(3):385--396, 1987.
\newblock doi:10.2969/jmsj/03930385.
\newblock Theorems 1-4 and 2-1, p.~388.

\bibitem{ShiodaRan}
T. Shioda, The Hodge conjecture for Fermat varieties,
\emph{Mathematische Annalen}, 245:175--184, 1979;
Z. Ran, Cycles on Fermat hypersurfaces,
\emph{Compositio Mathematica}, 42:121--142, 1981.
\newblock Concatenation and complementary-pair removal are stated in
Aoki, Theorem 1-4, p.~388.
% These original papers were not independently proof-audited here.

\bibitem{ShiodaFermat}
T. Shioda.
\newblock The Hodge conjecture for Fermat varieties.
\newblock \emph{Mathematische Annalen}, 245(2):175--184, 1979.
\newblock doi:10.1007/BF01428804.

\bibitem{ShiodaKatsura}
T. Shioda and T. Katsura.
\newblock On Fermat varieties.
\newblock \emph{Tohoku Mathematical Journal}, (2)31(1):97--115, 1979.
\newblock doi:10.2748/tmj/1178229881.

\bibitem{VanGeemenHalfTwist}
B. van Geemen.
\newblock Half twists of Hodge structures of CM-type.
\newblock \emph{Journal of the Mathematical Society of Japan}, 53(4):813--833, 2001.
\newblock doi:10.2969/jmsj/05340813.
\newblock \S1.4 and Proposition 2.8.

\bibitem{VillaflorPeriods}
R. Villaflor Loyola.
\newblock Periods of complete intersection algebraic cycles.
\newblock arXiv:1812.03964v5, 30 March 2021, Theorems 1.1--1.2.

\bibitem{KangFermat}
S.-J. Kang.
\newblock Refined motivic dimension of some Fermat varieties.
\newblock \emph{Bulletin of the Australian Mathematical Society}, 93(2):223--230, 2016.
\newblock doi:10.1017/S0004972715001379.
% Original full theorem/proof unavailable in this audit; Cor3.2 corroborated
% by Miranda et al., section3.4. Do not mark the original proof audited.

\bibitem{FermatLengths2026}
M. Miranda, H. Movasati, L. Rufino, and R. Villaflor.
\newblock Lengths of Hodge characters in Fermat varieties.
\newblock arXiv:2609.27301v2, 26 September 2026.
% Current preprint context; its integer programs were not replayed here.

\bibitem{MarkmanWeil}
E. Markman.
\newblock Cycles on abelian $2n$-folds of Weil type from secant sheaves on abelian $n$-folds.
\newblock arXiv:2502.03415v2, 8 June 2025.
\newblock Theorem 1.5.1 and Corollary 1.6.1.

\bibitem{MarkmanICM2026}
E. Markman.
\newblock Secant Sheaves and Weil Classes on Abelian Varieties.
\newblock \emph{Proceedings of the International Congress of Mathematicians 2026},
vol.~3, pp.~586--602.
\newblock doi:10.1137/25M1803796.
\newblock Theorem 1.2 and Corollary 1.3.

\bibitem{BrosnanDegeneration}
P. Brosnan.
\newblock Discriminants of Hermitian forms, maximal degenerations and Kontsevich's tropical approach to the Hodge conjecture.
\newblock arXiv:2609.14169v1, 12 September 2026.
% Exact Proposition4 unavailable in this audit; retain conditional scope.

\bibitem{SecantPreprintAudit}
D. Bhattacharjee and U. Bhattacharya.
\newblock Relative Secant Cycles and Hodge Classes.
\newblock \emph{Preprints.org}, 202602.0462, version~5, 5 August 2026.
\newblock \url{https://www.preprints.org/manuscript/202602.0462/v5}.
\newblock \S4.3, Proposition 2 (local eigencharacter identification audited).

\bibitem{SharifiDirichlet}
R. Sharifi.
\newblock Iwasawa Theory, Chapter 4: Cyclotomic fields.
\newblock Author's online notes, accessed 9 October 2026.
\newblock Lemma 4.1.18, Corollary 4.1.22, Proposition 4.2.13, Theorem 4.2.14.
\newblock \url{https://www.math.ucla.edu/~sharifi/notes/iwasawa-ch04.html}.

\bibitem{DLMFDirichlet}
NIST Digital Library of Mathematical Functions.
\newblock \S25.15, Dirichlet $L$-functions, equations 25.15.5 and 25.15.9.
\newblock \url{https://dlmf.nist.gov/25.15}, accessed 9 October 2026.
